\section{Conclusion and Future Work}
\label{sec:conclusion}

This paper implemented an Actual Usage Evictor as a global pre-eviction safety
gate for Kubernetes Descheduler. The plugin preserves the candidate ordering
and placement objective of the active strategy, but checks live CPU and memory
usage immediately before eviction. The implementation supports CPU and memory
ratios, while the present CPU-bound experiments empirically validate only the
CPU decision path. Request-relative thresholds distinguish currently active
Pods from idle candidates, while conservative handling of unavailable or
incomplete metrics prevents eviction without adequate runtime evidence.

In the evaluated single-replica API scenarios, each represented by one
independent run, the filter blocked the busy API and allowed idle candidates
to remain eligible. Controls produced failure rates of \(16.67\%\) and
\(17.24\%\) and observed failure spans of approximately \(4.8\,s\); both
treatments recorded no failed response. HNU retained one additional active
Node because no fallback existed on the protected source. LNU substituted an
idle Pod, completed the same three evictions, and retained the full \(85.7\%\)
memory-deviation reduction. For these CPU-bound workloads and tested layouts,
runtime filtering protected the original API instance while preserving
strategy progress when a compatible alternative candidate was available.

Future work should evaluate threshold sensitivity across differently
calibrated requests, incorporate temporal smoothing or multiple consecutive
samples into the filter itself, and study the trade-off between responsiveness
and stale observations. Evaluation should also cover replicated, stateful, and
production-derived workloads; heterogeneous Nodes; longer requests; and
additional Descheduler strategies. Finally, exposing block reasons and recent
usage ratios as metrics would improve operator visibility and support
workload-specific threshold tuning.
